\documentclass[UTF8,zihao=-4,openany]{ctexbook}

\usepackage[a4paper,margin=2.45cm,headheight=15pt]{geometry}
\usepackage{amsmath,amssymb,bm}
\usepackage{booktabs,tabularx,array}
\usepackage[dvipsnames]{xcolor}
\usepackage[most]{tcolorbox}
\usepackage{fancyhdr}
\usepackage{enumitem}
\usepackage{listings}
\usepackage[hidelinks]{hyperref}

\definecolor{Primary}{HTML}{355C7D}
\definecolor{Accent}{HTML}{6C8EAD}
\definecolor{SoftBlue}{HTML}{EDF4F8}
\definecolor{SoftGold}{HTML}{FFF8E8}
\definecolor{CodeBg}{HTML}{F6F7F8}

\setlength{\parindent}{2em}
\setlength{\parskip}{0.35em}
\setlist[itemize]{leftmargin=2em,itemsep=0.2em,topsep=0.3em}
\setlist[enumerate]{leftmargin=2.2em,itemsep=0.2em,topsep=0.3em}
\renewcommand{\arraystretch}{1.25}

\pagestyle{fancy}
\fancyhf{}
\fancyhead[L]{《动手学深度学习》PyTorch版}
\fancyhead[R]{提炼学习笔记（试作）}
\fancyfoot[C]{\thepage}

\hypersetup{
  pdftitle={《动手学深度学习》PyTorch版学习笔记},
  pdfauthor={learning-notes}
}

\lstset{
  basicstyle=\ttfamily\small,
  backgroundcolor=\color{CodeBg},
  frame=single,
  rulecolor=\color{Accent},
  breaklines=true,
  columns=fullflexible,
  keepspaces=true,
  showstringspaces=false,
  language=Python
}

\newtcolorbox{summarybox}[1][]{
  colback=SoftBlue,
  colframe=Primary,
  fonttitle=\bfseries,
  title={核心摘要},
  #1
}

\newtcolorbox{notebox}[1][]{
  colback=SoftGold,
  colframe=Goldenrod,
  fonttitle=\bfseries,
  title={易混淆点},
  #1
}

\newcommand{\code}[1]{\texttt{#1}}

\title{\heiti 《动手学深度学习》\\[0.4em]\Large PyTorch版学习笔记}
\author{内容提炼试作}
\date{\today}

\begin{document}

\maketitle
\clearpage
\pagenumbering{roman}

\chapter*{试作说明}
\addcontentsline{toc}{chapter}{试作说明}
本 PDF 用于检验笔记的内容密度、排版和编译流程，仅提炼原书第二章“预备知识”。
章内目录严格沿用 PyTorch 版原始顺序：
数据操作、数据预处理、线性代数、微积分、自动微分、概率、查阅文档。
正文只保留理解后续深度学习所需的概念、公式、PyTorch 操作和易混淆点。

\tableofcontents

\clearpage
\pagenumbering{arabic}
\setcounter{chapter}{1}
\chapter{预备知识}

本章把后续计算写成同一套对象与映射：数据是张量，前向是线性映射，训练目标是可微标量，
自动微分在计算图上实现链式法则，概率把预测写成分布。
每一公式先写如何由输入算出输出，再写这个数在问题里表示什么。

\section{数据操作}

记 $n$ 阶张量为 $\bm{X}\in\mathbb{R}^{d_1\times\cdots\times d_n}$。
$n=0,1,2$ 分别对应标量、向量、矩阵；更高阶用于批量、通道或时间轴。

\subsection{入门}

形状与元素总数为
\[
\operatorname{shape}(\bm{X})=(d_1,\ldots,d_n),
\qquad
\operatorname{numel}(\bm{X})=\prod_{i=1}^{n}d_i.
\]
$\operatorname{shape}$ 得到的是各轴长度，表示“有几条样本、几个特征、几条通道”等几何尺寸，
本身不是数据值。$\operatorname{numel}$ 得到的是一个正整数，表示这块内存里存了多少个数字。
\code{reshape} 只改变轴的划分，不改变元素值与存储次序。若
\[
\prod_{i=1}^{n}d_i=\prod_{j=1}^{m}d'_j,
\]
则存在把 $\mathbb{R}^{d_1\times\cdots\times d_n}$ 重新解释为 $\mathbb{R}^{d'_1\times\cdots\times d'_m}$ 的双射：
输出张量里每一个数仍是原来的测量值，只是被看成另一组轴上的条目。
某个 $d'_k=-1$ 时，由上式唯一确定该轴长度。

\begin{lstlisting}
x = torch.arange(12).reshape(3, 4)
x.shape, x.numel()
\end{lstlisting}
这里 $12$ 个数被排成 $3\times 4$：每一项仍是 $0,1,\ldots,11$ 中的那个整数，不是新算出来的量。

\subsection{运算符}

形状相同后，二元运算默认逐元素进行。若 $\bm{u},\bm{v}\in\mathbb{R}^{d}$，
\[
(f(\bm{u},\bm{v}))_i=f(u_i,v_i),\qquad i=1,\ldots,d.
\]
第 $i$ 个输出只由第 $i$ 个输入对算出。$u_i+v_i$ 是两个对应位置测量值之和；
$u_i v_i$ 是对应位置的乘积，仍是一个特征上的数，不是点积那种“整体相似度”。
沿轴 $k$ 连接要求除第 $k$ 轴外形状一致：
\[
\operatorname{cat}(\bm{A},\bm{B};k)
\in
\mathbb{R}^{d_1\times\cdots\times(d_k^{(A)}+d_k^{(B)})\times\cdots\times d_n}.
\]
连接不产生新数值，只把两段已有数据沿指定轴拼在一起。
比较返回 $\{0,1\}$：得到的 $1$ 表示“该位置关系成立”，$0$ 表示不成立。求和
\[
\sum(\bm{X})=\sum_{i_1=1}^{d_1}\cdots\sum_{i_n=1}^{d_n}X_{i_1\ldots i_n}
\]
把全部条目收成一个标量，表示这些数的总量（例如像素强度总和），不再保留位置。

\subsection{广播机制}

形状不同时，从末轴向前对齐。第 $i$ 轴可运算当且仅当
\[
d_i^{(X)}=d_i^{(Y)}
\quad\text{或}\quad
d_i^{(X)}=1
\quad\text{或}\quad
d_i^{(Y)}=1.
\]
长度为 $1$ 的轴在逻辑上复制到目标长度。于是
\[
\mathbb{R}^{3\times 1}+\mathbb{R}^{1\times 2}
\longrightarrow
\mathbb{R}^{3\times 2},
\qquad
(A+B)_{ij}=A_{i1}+B_{1j}.
\]
$(A+B)_{ij}$ 是“第 $i$ 个行标量”与“第 $j$ 个列标量”配对相加后的数；
几何上相当于把短轴上的值复制后逐元素相加，并不增加新的独立观测。

\subsection{索引和切片}

对 $\bm{X}\in\mathbb{R}^{m\times n}$，下标从 $0$ 起算。切片 $[a:b]$ 取 $\{a,a+1,\ldots,b-1\}$，
即半开区间 $[a,b)$，且 $-1$ 表示最后一个合法下标。多轴独立作用：
\[
\bm{X}[i,j]=X_{ij},\qquad
\bm{X}[a:b,:]\in\mathbb{R}^{(b-a)\times n}.
\]
$X_{ij}$ 就是第 $i$ 行第 $j$ 列那个原始测量值；切片得到的仍是原表中的一块，不是统计量。

\subsection{节省内存}

记 $\operatorname{id}(\cdot)$ 为对象标识。赋值
\[
\bm{Y}\leftarrow \bm{Y}+\bm{X}
\]
通常分配新存储后再改引用；切片写回
\[
\bm{Y}[:]\leftarrow \bm{Y}+\bm{X}
\]
保持 $\operatorname{id}(\bm{Y})$ 不变。这里比较的不是数值大小，而是“是否仍是同一块内存”。
$\operatorname{id}$ 相同表示后续原地修改会改到同一批数。

\subsection{转换为其他Python对象}

PyTorch 张量与 NumPy 数组可共享同一块内存，因此对一侧的原地更新会改变另一侧。
标量提取是映射
\[
\mathbb{R}^{1}\longrightarrow\mathbb{R},\qquad
\bm{x}\mapsto x=\code{item}(\bm{x}).
\]
得到的 $x$ 就是那一个张量条目的数值，供 Python 标量运算使用，不再带有梯度或形状信息。

\subsection{小结}

数据操作把结构写成 $\operatorname{shape}$、个数写成 $\operatorname{numel}$，
把组合写成逐元素映射与广播，把取值写成索引。
算出来的每一个条目，除非经过求和等归约，否则仍表示原来位置上的数据值。

\subsection{练习}

核验三件事：$\operatorname{numel}$ 这个正整数是否等于各轴长度之积、
广播后每个位置的数是哪两个输入配对相加、共享内存是否意味着改一侧会改另一侧的数。

\section{数据预处理}

原始表通常含缺失与非数值字段，不能直接作为上一节中的 $\bm{X}$。
预处理的目标是构造数值张量 $(\bm{X},\bm{y})$。

\subsection{读取数据集}

设表有 $n$ 行、$p$ 列。一行是一个样本，一列是一个属性。划分后
\[
\bm{x}^{(i)}\in\mathbb{R}^{p-1},\qquad y^{(i)}\in\mathbb{R},
\qquad i=1,\ldots,n.
\]
$\bm{x}^{(i)}$ 中的数是第 $i$ 个样本的输入特征；$y^{(i)}$ 是该样本要预测的目标值。
缺失在 pandas 中记为 $\mathrm{NaN}$，表示“这个位置没有观测”，不是数值 $0$。

\subsection{处理缺失值}

对数值列 $j$，令观测指标集 $\mathcal{I}_j=\{i:x_{ij}\text{ 非缺失}\}$。样本均值
\[
\bar{x}_j=\frac{1}{|\mathcal{I}_j|}\sum_{i\in\mathcal{I}_j}x_{ij}
\]
是该列已观测值的算术平均，表示“这一属性的典型水平”，不是某一次真实测量。
填充 $x_{ij}\leftarrow\bar{x}_j$ 只是用这个典型水平代替缺口。
类别列取值于有限集 $\mathcal{C}$ 时，独热映射为
\[
c\mapsto \bm{e}_c\in\{0,1\}^{|\mathcal{C}|},\qquad
(\bm{e}_c)_{c'}=[c=c'].
\]
向量里的 $1$ 表示“属于该类”，$0$ 表示“不属于”；各分量互斥，和为 $1$。
缺失可并入一个额外类别，使编码落在 $\mathbb{R}^{|\mathcal{C}|+1}$。

\subsection{转换为张量格式}

全部字段数值化且类型一致后，
\[
(\bm{x}^{(i)},y^{(i)})_{i=1}^{n}
\;\longrightarrow\;
\bm{X}\in\mathbb{R}^{n\times d},
\;
\bm{y}\in\mathbb{R}^{n}.
\]
$X_{ij}$ 是第 $i$ 个样本第 $j$ 个特征的数值；$y_i$ 是第 $i$ 个样本的标签。
二者已经是模型可做加减乘的实数，不再含字符串或 $\mathrm{NaN}$。

\subsection{小结}

预处理不改变学习目标 $y$ 的含义，只把不完整字段映成可计算的 $(\bm{X},\bm{y})$：
$\bm{X}$ 的条目是特征， $\bm{y}$ 的条目是要拟合的量。

\subsection{练习}

删除缺失最多的列是丢掉一维特征；均值填充则仍保留该维，但缺口处的数是列均值，不是观测。

\section{线性代数}

张量给出对象；线性代数给出映射。从标量升到矩阵，再由点积与乘法得到层间变换，
最后用范数把误差压成标量。

\subsection{标量}

标量 $x\in\mathbb{R}$ 满足
\[
x+y,\quad xy,\quad x/y\ (y\neq 0),\quad x^{p}.
\]
结果仍是一个实数：和、积、商、幂都表示在同一条数轴上的量，没有方向。

\subsection{向量}

$n$ 个标量排成列向量
\[
\bm{x}=
\begin{bmatrix}
x_1\\x_2\\\vdots\\x_n
\end{bmatrix}
\in\mathbb{R}^{n}.
\]
$x_i$ 是第 $i$ 个坐标上的分量，例如第 $i$ 个特征的取值；整个 $\bm{x}$ 表示一个有序的特征组合。

\subsubsection{长度、维度和形状}

$\bm{x}\in\mathbb{R}^{n}$ 表示分量数（向量维度）为 $n$，对应一维张量形状 $(n,)$。
这里的 $n$ 是“有几个特征”，不是向量的欧氏长度 $\|\bm{x}\|_2$。
张量的“维数”还可指轴数 $k$。

\subsection{矩阵}

\[
\bm{A}=
\begin{bmatrix}
a_{11}&\cdots&a_{1n}\\
\vdots&\ddots&\vdots\\
a_{m1}&\cdots&a_{mn}
\end{bmatrix}
\in\mathbb{R}^{m\times n}.
\]
$a_{ij}$ 是第 $i$ 行第 $j$ 列的数：在数据矩阵里常表示第 $i$ 个样本的第 $j$ 个特征；
在权重矩阵里常表示“输出单元 $i$ 对输入特征 $j$ 的权重”。
转置 $(\bm{A}^{\mathsf T})_{ij}=a_{ji}$ 只交换行与列的角色，不改变那些数本身。
若 $\bm{A}=\bm{A}^{\mathsf T}$，则 $\bm{A}$ 对称，此时 $m=n$。

\subsection{张量}

向量、矩阵分别是 $1$ 阶、$2$ 阶张量。图像 $\bm{X}\in\mathbb{R}^{H\times W\times C}$ 中，
每个 $X_{hwc}$ 是位置 $(h,w)$、通道 $c$ 上的像素强度（或特征响应）。
加入批量后 $\mathbb{R}^{B\times H\times W\times C}$ 的第一轴下标表示第几张图。

\subsection{张量算法的基本性质}

同形状逐元素运算保持形状。Hadamard 积
\[
(\bm{A}\odot\bm{B})_{ij}=a_{ij}b_{ij}
\]
得到的仍是位置 $(i,j)$ 上的一个数，表示两个对应条目的乘积，不是加权求和。
标量广播 $(\alpha+\bm{A})_{ij}=\alpha+a_{ij}$ 表示把同一个偏移加到每一个位置。

\subsection{降维}

对 $\bm{A}\in\mathbb{R}^{m\times n}$，
\[
\sum_{i=1}^{m}\sum_{j=1}^{n}a_{ij}
\]
是一个标量，表示全部条目的总和。沿轴归约
\[
\Bigl(\sum_{i=1}^{m}\bm{A}\Bigr)_{j}=\sum_{i=1}^{m}a_{ij}\in\mathbb{R}^{n},
\qquad
\Bigl(\sum_{j=1}^{n}\bm{A}\Bigr)_{i}=\sum_{j=1}^{n}a_{ij}\in\mathbb{R}^{m}
\]
分别表示“每一列的合计”和“每一行的合计”。均值
\[
\bar{A}=\frac{1}{mn}\sum_{i=1}^{m}\sum_{j=1}^{n}a_{ij}
\]
是这些数的平均水平，与单个条目同量纲。

\subsubsection{非降维求和}

行归一化
\[
\widetilde{a}_{ij}
=
\frac{a_{ij}}{\sum_{k=1}^{n}a_{ik}}
\quad
\Bigl(\sum_{k}a_{ik}\neq 0\Bigr)
\]
得到 $\widetilde{a}_{ij}\in[0,1]$（若原行为非负），且 $\sum_{j}\widetilde{a}_{ij}=1$。
它表示第 $i$ 行内部的相对比重，不再是原始量纲。累积和
\[
s_{ij}=\sum_{k=1}^{j}a_{ik}
\]
是该行从左到第 $j$ 列的部分和，表示“前 $j$ 项累计了多少”。

\subsection{点积（Dot Product）}

\[
\bm{x}^{\mathsf T}\bm{y}
=
\sum_{i=1}^{d}x_i y_i.
\]
结果是一个标量：把 $\bm{y}$ 的各分量用 $\bm{x}$ 作权重求和。
若 $\bm{x}$ 是权重、$\bm{y}$ 是特征，它就是线性层的一次打分；
若二者都是单位向量，该标量等于夹角余弦，取值 $[-1,1]$，表示方向有多一致。

\subsection{矩阵-向量积}

对 $\bm{A}\in\mathbb{R}^{m\times n}$、$\bm{x}\in\mathbb{R}^{n}$，
\[
(\bm{A}\bm{x})_i=\bm{a}_i^{\mathsf T}\bm{x}.
\]
第 $i$ 个输出是第 $i$ 个输出单元对输入 $\bm{x}$ 的线性响应（一个打分）。
整个 $\bm{A}\bm{x}\in\mathbb{R}^{m}$ 是 $m$ 个这样的打分。
线性性 $f(\alpha\bm{x}+\beta\bm{z})=\alpha f(\bm{x})+\beta f(\bm{z})$ 表示叠加输入则打分叠加。

\subsection{矩阵-矩阵乘法}

若 $\bm{A}\in\mathbb{R}^{n\times k}$、$\bm{B}\in\mathbb{R}^{k\times m}$，则
\[
c_{ij}=\bm{a}_i^{\mathsf T}\bm{b}_j=\sum_{\ell=1}^{k}a_{i\ell}b_{\ell j}.
\]
$c_{ij}$ 是 $\bm{A}$ 的第 $i$ 行与 $\bm{B}$ 的第 $j$ 列的点积：
一批输入同时通过同一线性层时，第 $i$ 个样本在第 $j$ 个输出通道上的响应。

\subsection{范数}

范数 $f:\mathbb{R}^{n}\to\mathbb{R}_{\ge 0}$ 满足绝对齐次、三角不等式与正定。
\[
\|\bm{x}\|_1=\sum_i|x_i|,
\qquad
\|\bm{x}\|_2=\sqrt{\sum_i x_i^2},
\qquad
\|\bm{X}\|_F=\sqrt{\sum_{i,j}x_{ij}^{2}}.
\]
这些数都是“有多大”：$\|\bm{x}\|_2$ 是欧氏长度，$\|\bm{x}\|_1$ 是各分量绝对值和（对异常值不那么敏感），
$\|\bm{X}\|_F$ 把矩阵当成向量后的欧氏长度。它们都 $\ge 0$，为零当且仅当对象是零。

\subsubsection{范数和目标}

\[
\min_{\bm{\theta}}
\bigl\|\bm{y}-f_{\bm{\theta}}(\bm{X})\bigr\|_2^2
\]
中的 $\|\bm{y}-\hat{\bm{y}}\|_2^2$ 是预测误差向量的平方欧氏长度，
数值越大表示整体预测离真值越远；它与单个 $y$ 同量纲的平方。

\subsection{关于线性代数的更多信息}

后续章节主要用到 $\bm{A}\bm{x}$ 的打分、$\bm{A}\bm{B}$ 的批量打分，以及 $\|\cdot\|$ 作为误差大小。

\subsection{小结}

线性代数把数据写成张量，把层写成点积/矩阵乘（输出是打分），把误差大小写成范数（非负标量）。

\subsection{练习}

核验恒等式两边每个条目是否同一数，以及结果形状是否等于“行点列”所要求的 $(n,m)$。

\section{微积分}

前向已写成映射 $f$；训练要知道参数扰动如何改变标量目标。

\subsection{导数和微分}

若极限存在，
\[
f'(x)
=
\lim_{h\to 0}\frac{f(x+h)-f(x)}{h}.
\]
$f'(x)$ 是一个实数：输入增加极小 $h$ 时，输出大约增加 $f'(x)\,h$。
正值表示局部上升，负值表示局部下降，绝对值是变化有多陡。
对 $f(x)=3x^2-4x$，
\[
f'(1)=2
\]
表示在 $x=1$ 处，输入增加 $\Delta x$ 则输出大约增加 $2\Delta x$。一阶展开
\[
f(x+\Delta x)=f(x)+f'(x)\Delta x+o(\Delta x)
\]
右边第二项就是用这个变化率作的线性近似。

\subsection{偏导数}

\[
\frac{\partial y}{\partial x_i}
=
\lim_{h\to 0}
\frac{
f(x_1,\ldots,x_i+h,\ldots,x_n)
-f(x_1,\ldots,x_i,\ldots,x_n)
}{h}.
\]
这是一个实数：只动第 $i$ 个自变量、其余固定时，$y$ 对 $x_i$ 的敏感度。
它不描述沿其他方向的变化。

\subsection{梯度}

\[
\nabla_{\bm{x}}f(\bm{x})
=
\begin{bmatrix}
\partial f/\partial x_1\\
\vdots\\
\partial f/\partial x_n
\end{bmatrix}.
\]
第 $i$ 个分量是 $f$ 对 $x_i$ 的偏导。整个向量指向 $f$ 上升最快的方向；
其长度表示沿该方向的变化有多快。一阶展开中
\[
\nabla_{\bm{x}}f(\bm{x})^{\mathsf T}\Delta\bm{x}
\]
就是“沿 $\Delta\bm{x}$ 走一步，$f$ 大约增加多少”。
$\nabla_{\bm{x}}\|\bm{x}\|_2^2=2\bm{x}$ 表示：平方欧氏长度对 $\bm{x}$ 的敏感度就是 $2\bm{x}$ 本身。

\subsection{链式法则}

\[
\frac{\mathrm{d}y}{\mathrm{d}x}
=
\frac{\mathrm{d}y}{\mathrm{d}u}
\frac{\mathrm{d}u}{\mathrm{d}x}.
\]
左边是复合后 $y$ 对最外层输入 $x$ 的总敏感度，等于“$y$ 对中间量 $u$ 的敏感度”乘“$u$ 对 $x$ 的敏感度”。
多路径时
\[
\frac{\partial y}{\partial x_i}
=
\sum_{j=1}^{m}
\frac{\partial y}{\partial u_j}
\frac{\partial u_j}{\partial x_i}
\]
把每条路径的敏感度相加。反向传播算出的就是这个总敏感度，用来决定参数该往哪边改。

\subsection{小结}

$f'(x)$ 是切线斜率（输出随输入变多快）；偏导是沿单一坐标的斜率；
梯度把这些斜率收成一个向量；链式法则把各层斜率连成对参数的总斜率。

\subsection{练习}

画出 $f$ 与切线时，切线斜率必须等于所求的 $f'(x)$ 那个数。

\section{自动微分}

前向记录计算图 $\mathcal{G}$，反向对每个节点应用
\[
\frac{\partial L}{\partial u}
=
\sum_{v:\,u\to v}
\frac{\partial L}{\partial v}
\frac{\partial v}{\partial u}.
\]
左边是损失 $L$ 对中间量 $u$ 的敏感度，由下游敏感度与局部导数相乘再相加得到。

\subsection{一个简单的例子}

\[
y=2\bm{x}^{\mathsf T}\bm{x}=2\sum_i x_i^2,
\qquad
\nabla_{\bm{x}}y=4\bm{x}.
\]
$y$ 是一个非负标量，表示 $2\|\bm{x}\|_2^2$。$\nabla_{\bm{x}}y$ 的第 $i$ 个分量 $4x_i$ 表示：
$x_i$ 增加 $1$，$y$ 大约增加 $4x_i$。\code{x.grad} 读到的就是这些数。
默认累积：第二次反向若不清零，\code{x.grad} 变成 $8\bm{x}$，那是两次敏感度相加，不再是单次导数。

\subsection{非标量变量的反向传播}

Jacobian $J_{ij}=\partial y_i/\partial x_j$ 表示“第 $i$ 个输出对第 $j$ 个输入的敏感度”。
训练通常算 $\bm{v}^{\mathsf T}J$：得到的向量是“把各输出用 $\bm{v}$ 加权后的总损失”对 $\bm{x}$ 的梯度。
取 $\bm{v}=\bm{1}$ 时，
\[
\nabla_{\bm{x}}\sum_i x_i^2=2\bm{x}
\]
的含义是：把所有输出分量同等看待后加总，再问 $\bm{x}$ 往哪边能减小这个总和。

\subsection{分离计算}

$\bm{u}=\mathrm{detach}(\bm{y})$ 后 $\nabla_{\bm{x}}z=\bm{u}$。
这里每个分量 $u_i$ 被当作常数权重，梯度不再包含“$\bm{u}$ 本身还依赖 $\bm{x}$”那一条路。
因此得到的数是部分导数，不是完整复合 $z=\sum_i x_i^3$ 的 $3x_i^2$。

\subsection{Python控制流的梯度计算}

若本次路径上 $f(a)=ka$，则
\[
\frac{\mathrm{d}f}{\mathrm{d}a}=k=\frac{f(a)}{a},\qquad a\neq 0.
\]
$k$ 这个数就是本次实际执行的那段仿射关系的斜率；换一条分支，$k$ 会变。

\subsection{小结}

\code{x.grad} 里的每一个数都是 $\partial L/\partial x_i$：
对应参数增大时，标量损失大约增加多少。要减小损失，应沿其负方向更新。

\subsection{练习}

核验二阶时每个海森条目 $\partial^2 L/\partial x_i\partial x_j$ 表示一阶敏感度再对 $x_j$ 的变化；
重复反向时 \code{grad} 中的数会叠加。

\subsection{练习解答}

\subsubsection{1. 为什么二阶导数的开销更大？}

海森 $H_{ij}$ 有最多 $n^2$ 个数，每个都是“梯度的第 $i$ 个分量再对 $x_j$ 求导”。
要保存并再遍历一阶图，时间和内存都按这一矩阵规模增加。

\subsubsection{2. 连续运行两次反向传播}

第二次得到的 \code{grad} 是两次 $\nabla L$ 之和，数值上约为单次的两倍，不是新的一次导数。

\subsubsection{3. 将控制流例子中的输入改成向量或矩阵}

$\bm{1}^{\mathsf T}\bm{f}(\bm{a})$ 把各输出加总成一个标量；其梯度每个分量仍是该路径上的斜率 $k$，
含义不变：对应输入增大时，输出总和大约增加 $k$。

\subsubsection{4. 重新设计控制流梯度例子}

$x>0$ 时 $g'(x)=2x$ 表示抛物线的切线斜率；$x<0$ 时 $-3x^2$ 是另一支的斜率。
分支选择本身没有导数。

\section{概率}

点估计给出 $\hat{y}$；概率把不确定性写成 $[0,1]$ 上的数。

\subsection{基本概率论}

\[
\widehat{P}_N(\mathcal{A})
=
\frac{1}{N}\sum_{i=1}^{N}\mathbf{1}\{X_i\in\mathcal{A}\}
\]
是 $N$ 次独立试验中事件发生的频率，取值 $[0,1]$。大数定律下它逼近真实概率 $P(\mathcal{A})$，
即“该事件长期发生的比例”。

\subsubsection{概率论公理}

$P(\mathcal{A})\in[0,1]$，$P(\mathcal{S})=1$ 表示“必然发生”。
互斥并的概率相加，表示把不重叠情形的可能性拼起来。$P(\varnothing)=0$ 表示不可能事件。

\subsubsection{随机变量}

$P(X=x)$ 是离散取值恰好为 $x$ 的概率。连续时
\[
P(a\le X\le b)=\int_{a}^{b}p(x)\,\mathrm{d}x
\]
是落在区间内的概率；$p(x)$ 本身是密度，可以大于 $1$，不是概率。

\subsection{处理多个随机变量}

联合、条件、边际都是 $[0,1]$ 上的数，但条件不同。

\subsubsection{联合概率}

$P(A=a,B=b)$ 是两件事同时发生的概率，不会超过任一边际。

\subsubsection{条件概率}

\[
P(B=b\mid A=a)
=
\frac{P(A=a,B=b)}{P(A=a)}.
\]
这是在已经知道 $A=a$ 的子总体里，$B=b$ 所占的比例，不是无条件下的发生频率。

\subsubsection{贝叶斯定理}

\[
P(A\mid B)
=
\frac{P(B\mid A)P(A)}{P(B)}.
\]
$P(A)$ 是见到 $B$ 之前认为 $A$ 为真的程度；$P(B\mid A)$ 是 $A$ 为真时看到 $B$ 的可能性；
$P(A\mid B)$ 是见到 $B$ 之后更新过的相信程度，仍在 $[0,1]$。

\subsubsection{边际化}

$P(B)=\sum_A P(A,B)$ 把 $A$ 的各种可能加总，得到“不管 $A$ 如何，$B$ 发生的总概率”。

\subsubsection{独立性}

$P(A,B)=P(A)P(B)$ 表示联合恰好是边际的乘积：知道 $B$ 不会改变对 $A$ 的概率，
即 $P(A\mid B)=P(A)$ 这个数保持不变。

\subsubsection{应用}

$P(H=1)=0.0015$ 表示感染的先验比例约 $0.15\%$。
$P(D_1=1)=0.011485$ 表示一次检测呈阳性的总比例。
后验
\[
P(H=1\mid D_1=1)\approx 0.1306
\]
表示：在已经阳性的人里，真正感染的比例约 $13\%$，不是 $100\%$。
灵敏度高只说明“有病时几乎一定阳性”，不能把阳性直接读成“一定有病”。

\subsection{期望和方差}

\[
\mathbb{E}[X]=\sum_x xP(X=x)
\]
是按概率加权的平均水平，与 $X$ 同量纲。
\[
\operatorname{Var}(X)=\mathbb{E}[(X-\mathbb{E}[X])^2]
\]
是相对均值的平方偏离的平均，量纲是 $X$ 的平方；标准差开方后与 $X$ 同量纲，表示典型起伏有多大。

\subsection{小结}

概率值是 $[0,1]$ 上的可能性；条件概率是限制子总体后的比例；后验是用数据更新后的相信程度；
期望是平均水平，方差是分散程度。

\subsection{练习}

核验频率 $\widehat{P}_N$ 是否逼近概率，以及后验是否真的被先验和似然共同决定。

\section{查阅文档}

接口随版本变化，核验对象本身。这里没有新的数值，只核验函数签名是否对应前面的数学映射。

\subsection{查找模块中的所有函数和类}

\code{dir(module)} 返回名称列表，不是数值结果。

\subsection{查找特定函数和类的用法}

\code{help} 给出参数与返回值说明，用来核对实现是否就是本节公式所定义的那个映射。

\subsection{小结}

先列出候选对象，再读签名，确认输入输出的数学含义与文档一致。

\subsection{练习}

对选定张量函数：文档中的返回值对应公式里的哪一个数（形状、归约轴、是否就地修改）。


\end{document}
